\section{Lecture 1: Randomized Controlled Trials}
\sourcepages{Lecture1-10.pdf, pp. 2, 1, 3; reading order: main text on p. 2, its supplement on p. 1, then p. 3.}
Regression in an RCT analysis of the average treatment effect reduces variance; it does not play a role in identifying the ATE.
\[
\Delta_i=Y_i(1)-Y_i(0),\qquad \tau=E[Y_i(1)-Y_i(0)].
\]
Estimate and construct inference for $\tau$ using the difference-in-means estimator:
\[
\widehat\tau_{\rm DM}=\frac1{n_1}\sum_{W_i=1}Y_i-\frac1{n_0}\sum_{W_i=0}Y_i,
\qquad n_w=|\{i:W_i=w\}|.
\]
The note marks this estimator as consistent and unbiased.
\paragraph{Supplement: assumptions.}
\[
Y_i=Y_i(W_i),\qquad W_i\perp\{Y_i(0),Y_i(1)\}.
\]
The first is labeled SUTVA; the second is random treatment assignment. SUTVA means the stable unit treatment value assumption: (1) unit stability, with experimental units independent of one another; (2) a single treatment level/version. These are the author's short descriptions.
\subsection{Linear outcome model and variance reduction}
\[
Y_i(w)=c_{(w)}+X_i\beta_{(w)}+\varepsilon_i(w),\quad
E[\varepsilon_i(w)\mid X_i]=0,\quad
\operatorname{Var}[\varepsilon_i(w)\mid X_i]=b^2.
\]
The coefficients are fitted separately for the two treatment states. Write $A=\operatorname{Var}[X]$ and $\|v\|_A^2=v'Av$. The variance in the notes is
\[
V_{\rm DM}=4b^2+\|\beta_{(0)}+\beta_{(1)}\|_A^2
 +\|\beta_{(0)}-\beta_{(1)}\|_A^2.
\]
Also,
\[
\tau=c_{(1)}-c_{(0)}+E[X](\beta_{(1)}-\beta_{(0)}),\qquad
\widehat\tau_{\rm OLS}=\widehat c_{(1)}-\widehat c_{(0)}
 +\overline X(\widehat\beta_{(1)}-\widehat\beta_{(0)}).
\]
The error decomposition written in the notebook is
\begin{align*}
\widehat\tau_{\rm OLS}-\tau={}&(\widehat c_{(1)}-c_{(1)})-(\widehat c_{(0)}-c_{(0)})
 +\overline X(\beta_{(1)}-\beta_{(0)})\\
&+\overline X(\widehat\beta_{(1)}-\widehat\beta_{(0)}-\beta_{(1)}+\beta_{(0)}).
\end{align*}
The first two terms are annotated $N(0,b^2/n_1)$ and $N(0,b^2/n_0)$; the third is annotated $N(0,\|\beta_{(1)}-\beta_{(0)}\|_A^2/n)$; the last is $O_p(1/n)$. The margin says that uncertainty in the estimated slopes tends to zero. Thus
\[
\sqrt n(\widehat\tau_{\rm OLS}-\tau)\Rightarrow N(0,V_{\rm OLS}),\qquad
V_{\rm OLS}=4b^2+\|\beta_{(0)}-\beta_{(1)}\|_A^2,
\]
\[
V_{\rm DM}=V_{\rm OLS}+\|\beta_{(0)}+\beta_{(1)}\|_A^2.
\]
\textit{Translator's note: the displayed error decomposition and intercept statements use centered covariates, although this centering is not explicitly restated on the scanned page. The factors of four correspond to the balanced-allocation calculation in the notes.}
\subsection{OLS under a general outcome regression}
The next page says that OLS is better than difference in means in this model. Suppose
\[
Y_i(w)=\mu_{(w)}(X_i)+\varepsilon_i(w),\quad
E[\varepsilon_i(w)\mid X_i]=0,\quad
\operatorname{Var}[\varepsilon_i(w)\mid X_i]=b^2.
\]
For difference in means,
\[
\sqrt n(\widehat\tau_{\rm DM}-\tau)\Rightarrow N(0,V_{\rm DM}),\quad
V_{\rm DM}=4b^2+2\operatorname{Var}[\mu_{(0)}(X_i)]
+2\operatorname{Var}[\mu_{(1)}(X_i)].
\]
For OLS, the Huber--White calculation starts from the population projection
\[
(c^*_{(w)},\beta^*_{(w)})=arg\min_{c,\beta}
 E[(Y_i(w)-X_i\beta-c)^2].
\]
The notebook records
\[
\sqrt{n_w}\left\{\begin{pmatrix}\widehat c_{(w)}\\\widehat\beta_{(w)}\end{pmatrix}
-\begin{pmatrix}c^*_{(w)}\\\beta^*_{(w)}\end{pmatrix}\right\}
\Rightarrow N\left(0,\begin{pmatrix}\operatorname{MSE}^*_{(w)}&0\\0&\cdots\end{pmatrix}\right),
\]
with $c^*_{(w)}=E[Y_i(w)]$ and
\[
\operatorname{MSE}^*_{(w)}=E[(Y_i(w)-X_i\beta^*_{(w)}-c^*_{(w)})^2].
\]
It concludes
\[
\sqrt n(\widehat\tau_{\rm OLS}-\tau)\Rightarrow N(0,V_{\rm OLS}),\qquad
V_{\rm OLS}=V_{\rm DM}-\|\beta^*_{(0)}+\beta^*_{(1)}\|_A^2.
\]
The red annotation says ``$V_{\rm OLS}<V_{\rm DM}$ always holds.''
\textit{Translator's note: the displayed squared-norm identity gives a weak inequality, with equality when the norm is zero. The blanket strict inequality and the simplified covariance matrix above are preserved as statements in the original notes.}
\paragraph{Personal observation.}
OLS can be viewed as
\[
\widehat\tau_{\rm OLS}=\frac1n\sum_{i=1}^n
\{(\widehat c_{(1)}+X_i\widehat\beta_{(1)})
 -(\widehat c_{(0)}+X_i\widehat\beta_{(0)})\}
=\frac1n\sum_i[\widehat\mu_{(1)}(X_i)-\widehat\mu_{(0)}(X_i)].
\]
Here $\widehat\mu_{(w)}(x)$ is the fitted OLS value at $x$. It could be replaced with another method.

\section{Lecture 2: Unconfoundedness and the Propensity Score}
\sourcepages{Lecture1-10.pdf, pp. 5, 4, 7, 6; p. 4 continues the bottom of p. 5, and p. 6 supplements p. 7.}
Unconfoundedness: the aim is to remove the influence of confounding factors. Propensity score: the probability of treatment given the covariates.
\subsection{Aggregating difference-in-means estimators}
The red margin describes this as simply weighting by probability. With discrete, numerically valued $X_i$, assume
\[
\{Y_i(0),Y_i(1)\}\perp W_i\mid X_i=x,\qquad x\in\mathcal X.
\]
Here $x$ is the grouping variable. Define the group-wise ATE
\[
\tau(x)=E[Y_i(1)-Y_i(0)\mid X_i=x].
\]
Then
\begin{align*}
\widehat\tau_{\rm AGG}&=\sum_{x\in\mathcal X}\frac{n_x}{n}\widehat\tau(x),\\
\widehat\tau(x)&=\frac1{n_{x1}}\sum_{X_i=x,W_i=1}Y_i
-\frac1{n_{x0}}\sum_{X_i=x,W_i=0}Y_i,\\
n_x&=|\{i:X_i=x\}|,\qquad
n_{xw}=|\{i:X_i=x,W_i=w\}|.
\end{align*}
Let $e(x)=P(W_i=1\mid X_i=x)$. The within-group limit is
\begin{align*}
\sqrt{n_x}(\widehat\tau(x)-\tau(x))
&\Rightarrow N\left(0,
\frac{\operatorname{Var}[Y_i(0)\mid X_i=x]}{1-e(x)}
+\frac{\operatorname{Var}[Y_i(1)\mid X_i=x]}{e(x)}\right)\\
&=N\left(0,\frac{b^2(x)}{e(x)(1-e(x))}\right),
\end{align*}
where the second equality uses the common conditional variance written in the notes. Set $\pi(x)=P(X_i=x)$ and $\widehat\pi(x)=n_x/n$. Then
\begin{align*}
\widehat\tau_{\rm AGG}={}&\sum_x\pi(x)\tau(x)
+\sum_x\pi(x)(\widehat\tau(x)-\tau(x))\\
&+\sum_x(\widehat\pi(x)-\pi(x))\tau(x)
+\sum_x(\widehat\pi(x)-\pi(x))(\widehat\tau(x)-\tau(x)).
\end{align*}
The first term is $\tau$. The second is annotated as asymptotically normal, with variance
\[
\sum_x\pi^2(x)\operatorname{Var}[\widehat\tau(x)].
\]
The third has variance $n^{-1}\operatorname{Var}[\tau(X_i)]$; the last is $O_p(1/n)$. Following the arrow onto the left page,
\[
\sqrt n(\widehat\tau_{\rm AGG}-\tau)\Rightarrow N(0,V_{\rm AGG}),
\]
\begin{align*}
V_{\rm AGG}&=\operatorname{Var}[\tau(X_i)]
+\sum_{x\in\mathcal X}\pi^2(x)\frac1{\pi(x)}
\frac{b^2(x)}{e(x)(1-e(x))}\\
&=\operatorname{Var}[\tau(X_i)]
+E\left[\frac{b^2(X_i)}{e(X_i)(1-e(X_i))}\right].
\end{align*}
\subsection{Continuous covariates and propensity scores}
Given $X_i$, $W_i$ cannot ``peek'' at $Y_i$: this is the author's intuition for unconfoundedness. If
\[
\{Y_i(0),Y_i(1)\}\perp W_i\mid X_i,
\quad e(x)=P(W_i=1\mid X_i=x),
\]
then
\[
\{Y_i(0),Y_i(1)\}\perp W_i\mid e(X_i).
\]
\paragraph{Supplementary sketches and calculation.}
For discrete $X$, the sketch shows two outcome points for each $x_1,x_2,x_3$, joined by a downward arrow. The upper values are 10, 8, 7, the lower values 5, 4, 5, and the gaps are 5, 4, 2. For continuous $X$, the two sets of points become upper and lower curves; comparisons are made over $x\in\mathcal X$ and related to $e(x)$.
The proof written beneath the sketches is
\begin{align*}
&P(W_i=w\mid\{Y_i(0),Y_i(1)\},e(X_i))\\
&\quad=\int_{\mathcal X}P(W_i=w\mid\{Y_i(0),Y_i(1)\},X_i=x)
 P(X_i=x\mid e(X_i))\,dx\\
&\quad=\int_{\mathcal X}P(W_i=w\mid X_i=x)P(X_i=x\mid e(X_i))\,dx\\
&\quad=e(X_i)\mathbf1_{\{w=1\}}+(1-e(X_i))\mathbf1_{\{w=0\}}.
\end{align*}
The margin labels this ``Lebesgue integration.''
\textit{Translator's note: this reproduces the notebook's informal integral notation; the first conditional mixing distribution omits conditioning on the potential outcomes. The final expression is the intended balancing-score result.}
\subsection{Propensity stratification}
\begin{enumerate}
\item Sort the observations so that $\widehat e(X_{i_1})\leq\widehat e(X_{i_2})\leq\cdots\leq\widehat e(X_{i_n})$.
\item Split them into $J$ strata and estimate a difference in means within each stratum. If $I_j$ contains ranks $\lfloor(j-1)n/J\rfloor+1$ through $\lfloor jn/J\rfloor$, write the abbreviated formula on the page as
\[
\widehat\tau_j=\frac{\sum_{i\in I_j}W_iY_i}{\sum_{i\in I_j}W_i}
-\frac{\sum_{i\in I_j}(1-W_i)Y_i}{\sum_{i\in I_j}(1-W_i)}.
\]
The blue comment says that stratification makes the control and treatment groups more similar.
\item Average the stratum estimates: $\widehat\tau_{\rm STRAT}=J^{-1}\sum_{j=1}^J\widehat\tau_j$. The notebook calls this consistent.
\end{enumerate}
\subsection{Inverse-propensity weighting}
\sourcepages{Lecture1-10.pdf, pp. 9, 8, 10; p. 8 supplies the model used in the variance calculation on p. 9.}
\[
\widehat\tau_{\rm IPW}=\frac1n\sum_{i=1}^n
\left\{\frac{W_iY_i}{\widehat e(X_i)}-\frac{(1-W_i)Y_i}{1-\widehat e(X_i)}\right\}.
\]
Oracle IPW assumes $e(x)$ is perfectly known:
\[
\widehat\tau^*_{\rm IPW}=\frac1n\sum_i
\left\{\frac{W_iY_i}{e(X_i)}-\frac{(1-W_i)Y_i}{1-e(X_i)}\right\}.
\]
Suppose overlap holds, $\eta\leq e(x)\leq1-\eta$ for all $x\in\mathcal X$, and $|Y_i|<M$. If
\[
\sup_{x\in\mathcal X}|e(x)-\widehat e(x)|=O_p(a_n)\longrightarrow0,
\]
the notebook gives
\[
|\widehat\tau_{\rm IPW}-\widehat\tau^*_{\rm IPW}|=O_p(a_nM/\eta).
\]
Thus consistency of the oracle estimator transfers to estimated-propensity IPW. Moreover,
\[
E[\widehat\tau^*_{\rm IPW}]=E[Y_i(1)-Y_i(0)]=\tau.
\]
With overlap, oracle IPW converges at the $1/\sqrt n$ scale.
\paragraph{Variance calculation: supplementary model.}
The left page supplies
\begin{align*}
Y_i(0)&=c(X_i)-(1-e(X_i))\tau(X_i)+\varepsilon_i(0),\\
Y_i(1)&=c(X_i)+e(X_i)\tau(X_i)+\varepsilon_i(1),\\
E[\varepsilon_i(w)\mid X_i]&=0,\qquad
\operatorname{Var}[\varepsilon_i(w)\mid X_i=x]=b^2(x).
\end{align*}
Then the notes state
\begin{align*}
n\operatorname{Var}[\widehat\tau^*_{\rm IPW}]
=V_{\rm IPW^*}
={}&E\left[\frac{c^2(X_i)}{e(X_i)(1-e(X_i))}\right]
+\operatorname{Var}[\tau(X_i)]\\
&+E\left[\frac{b^2(X_i)}{e(X_i)(1-e(X_i))}\right],\\
\sqrt n(\widehat\tau^*_{\rm IPW}-\tau)&\Rightarrow N(0,V_{\rm IPW^*}).
\end{align*}
\paragraph{Personal comparison.}
$c(x)$ describes the relationship between $X$ and $Y$, rather than the relationship between $X$ and $\Delta Y$. Compared with aggregation,
\[
V_{\rm IPW^*}=V_{\rm AGG}+E\left[\frac{c^2(X_i)}{e(X_i)(1-e(X_i))}\right].
\]
Unless $c(X_i)$ is zero everywhere, oracle IPW has worse variance than aggregation. Compared with linear modeling: IPW is based on $\{Y_i(0),Y_i(1)\}\perp W_i\mid X_i$, while the linear model relies on linearity.
\textit{Translator's note: the two displayed potential-outcome equations and the claimed $c^2$ variance expression are transcribed as written. The weighting defining $c(x)$ needs checking: the $c^2$ expression ordinarily uses $c(x)=(1-e(x))\mu_1(x)+e(x)\mu_0(x)$, whereas the scanned equations use the opposite coefficients.}

\section{Lecture 3: Efficient Treatment Effect Estimation via Augmented IPW}
\sourcepages{Lecture1-10.pdf, pp. 11--12.}
Ordinary IPW has poor large-sample properties. For example, with discrete $X$, it is worse than the aggregate estimator $\widehat\tau$.
\subsection{Two characterizations of the ATE}
\begin{enumerate}
\item Using $e(x)=P(W_i=1\mid X_i=x)$, $\tau=E[\widehat\tau^*_{\rm IPW}]$, where the oracle estimator is as above.
\item Using $\mu_{(w)}(x)=E[Y_i(w)\mid X_i=x]$,
\[
\tau(x)=\mu_{(1)}(x)-\mu_{(0)}(x),\quad
\tau=E[\mu_{(1)}(X_i)-\mu_{(0)}(X_i)],
\]
which motivates
\[
\widehat\tau_{\rm REG}=\frac1n\sum_i[\widehat\mu_{(1)}(X_i)-\widehat\mu_{(0)}(X_i)].
\]
\end{enumerate}
\subsection{The augmented IPW estimator}
\begin{align*}
\widehat\tau_{\rm AIPW}=\frac1n\sum_i\bigg[&\widehat\mu_{(1)}(X_i)-\widehat\mu_{(0)}(X_i)
+W_i\frac{Y_i-\widehat\mu_{(1)}(X_i)}{\widehat e(X_i)}\\
&-(1-W_i)\frac{Y_i-\widehat\mu_{(0)}(X_i)}{1-\widehat e(X_i)}\bigg].
\end{align*}
\paragraph{Personal intuition.}
First use $\mu(x)$ to estimate the outcome, and then use IPW on the regression residuals. IPW makes the covariate composition of the data ``fair.''
\subsection{Double robustness}
AIPW is consistent if either the outcome regressions $\widehat\mu_{(w)}(x)$ are consistent or $\widehat e(x)$ is consistent.
When $\widehat\mu_{(w)}(x)\approx\mu_{(w)}(x)$, the regression average is consistent and the weighted residual terms have mean approximately zero because
\[
E[Y_i-\widehat\mu_{(W_i)}(X_i)\mid X_i,W_i]\approx0.
\]
The propensity term is held fixed by this conditioning. When instead $\widehat e(x)\approx e(x)$, rearrange:
\begin{align*}
\widehat\tau_{\rm AIPW}={}&\frac1n\sum_i\left[\frac{W_iY_i}{\widehat e(X_i)}
-\frac{(1-W_i)Y_i}{1-\widehat e(X_i)}\right]\\
&+\frac1n\sum_i\left[\widehat\mu_{(1)}(X_i)\left(1-\frac{W_i}{\widehat e(X_i)}\right)
-\widehat\mu_{(0)}(X_i)\left(1-\frac{1-W_i}{1-\widehat e(X_i)}\right)\right].
\end{align*}
The first part is IPW. The second has mean approximately zero since
$E[1-W_i/\widehat e(X_i)\mid X_i]\approx0$; conditioning holds the fitted outcome regression fixed.
\subsection{Semiparametric efficiency}
The notes call AIPW asymptotically optimal among nonparametric estimators. When the nuisance functions are estimated sufficiently accurately,
\[
\sqrt n(\widehat\tau_{\rm AIPW}-\widehat\tau^*_{\rm AIPW})\longrightarrow_p0,
\]
where the starred estimator is oracle AIPW. This oracle estimator is just an IID average, so the central limit theorem also applies to the estimated version:
\[
\sqrt n(\widehat\tau_{\rm AIPW}-\tau)\Rightarrow N(0,V^*),
\]
\[
V^*=\operatorname{Var}[\tau(X_i)]
+E\left[\frac{b_0^2(X_i)}{1-e(X_i)}\right]
+E\left[\frac{b_1^2(X_i)}{e(X_i)}\right].
\]
The margin describes this as a Cramer--Rao-type constraint: this is the best attainable effect/variance.
\subsection{AIPW and cross-fitting}
\sourcepages{Lecture1-10.pdf, pp. 13--14.}
Randomly split the data into two halves, $I_1$ and $I_2$. The annotation says to split the data and choose the halves at random. Set
\[
\widehat\tau_{\rm AIPW}=\frac{|I_1|}{n}\widehat\tau^{I_1}
+\frac{|I_2|}{n}\widehat\tau^{I_2},
\]
\begin{align*}
\widehat\tau^{I_1}=\frac1{|I_1|}\sum_{i\in I_1}\bigg[&\widehat\mu_{(1)}^{I_2}(X_i)-\widehat\mu_{(0)}^{I_2}(X_i)
+W_i\frac{Y_i-\widehat\mu_{(1)}^{I_2}(X_i)}{\widehat e^{I_2}(X_i)}\\
&-(1-W_i)\frac{Y_i-\widehat\mu_{(0)}^{I_2}(X_i)}{1-\widehat e^{I_2}(X_i)}\bigg].
\end{align*}
Asymptotically, we can ignore the specific properties of the chosen machine-learning adjustment method and instead rely on:
\begin{enumerate}
\item Overlap: $\eta<e(x)<1-\eta$ for all $x\in\mathcal X$, with $\eta>0$.
\item Consistency of the nuisance estimators:
\[
\sup_x|\widehat\mu_{(w)}^{I_2}(x)-\mu_{(w)}(x)|\longrightarrow_p0,
\quad \sup_x|\widehat e^{I_2}(x)-e(x)|\longrightarrow_p0.
\]
The blue sentence says that the machine-learning adjustment methods have large-sample consistency.
\item Risk decay for the outcome and propensity models:
\[
E[(\widehat\mu_{(w)}^{I_2}(X_i)-\mu_{(w)}(X_i))^2]\,
E[(\widehat e^{I_2}(X_i)-e(X_i))^2]=o(1/n).
\]
If both are $\sqrt n$-consistent, the product is of order $1/n^2$. A simple sufficient approach is for both regression adjustments to converge at $o(n^{-1/4})$ in RMSE.
\end{enumerate}
\paragraph{Proof recorded on the continuation page.}
Write the oracle estimator as
$\widehat\tau^*=|I_1|\widehat\tau^{I_1,*}/n+|I_2|\widehat\tau^{I_2,*}/n$.
To establish $\sqrt n(\widehat\tau_{\rm AIPW}-\widehat\tau^*)\to_p0$, it suffices to show the corresponding result for the treated mean; the other components are analogous. Define
\[
\widehat\tau^{I_1}=\widehat\mu^{I_1}_{(1)}-\widehat\mu^{I_1}_{(0)},\qquad
\widehat\mu^{I_1}_{(1)}=\frac1{|I_1|}\sum_{i\in I_1}
\left[\widehat\mu^{I_2}_{(1)}(X_i)+W_i\frac{Y_i-\widehat\mu^{I_2}_{(1)}(X_i)}{\widehat e^{I_2}(X_i)}\right].
\]
Subtracting its oracle version gives
\begin{align*}
\widehat\mu^{I_1}_{(1)}-\widehat\mu^{I_1,*}_{(1)}
={}&\frac1{|I_1|}\sum_{i\in I_1}
(\widehat\mu^{I_2}_{(1)}(X_i)-\mu_{(1)}(X_i))\left(1-\frac{W_i}{e(X_i)}\right)\\
&+\frac1{|I_1|}\sum_{i\in I_1}W_i(Y_i-\mu_{(1)}(X_i))
\left(\frac1{\widehat e^{I_2}(X_i)}-\frac1{e(X_i)}\right)\\
&-\frac1{|I_1|}\sum_{i\in I_1}W_i(\widehat\mu^{I_2}_{(1)}(X_i)-\mu_{(1)}(X_i))
\left(\frac1{\widehat e^{I_2}(X_i)}-\frac1{e(X_i)}\right).
\end{align*}
The blue annotations bound the first contribution by an expression written as $o_p(1)/n$ and the last by $o_p(1/\sqrt n)$.
\textit{Translator's note: it is unclear whether the first handwritten bound is intended for the term or its variance. The target statement is the $o_p(n^{-1/2})$ remainder stated above.}
For $K$ folds, use nuisance estimates trained without the fold $k(i)$ containing observation $i$:
\begin{align*}
\widehat\tau_{\rm AIPW}=\frac1n\sum_i\bigg[&\widehat\mu_{(1)}^{(-k(i))}(X_i)-\widehat\mu_{(0)}^{(-k(i))}(X_i)
+W_i\frac{Y_i-\widehat\mu_{(1)}^{(-k(i))}(X_i)}{\widehat e^{(-k(i))}(X_i)}\\
&-(1-W_i)\frac{Y_i-\widehat\mu_{(0)}^{(-k(i))}(X_i)}{1-\widehat e^{(-k(i))}(X_i)}\bigg].
\end{align*}
\subsection{Confidence intervals}
\sourcepages{Lecture1-10.pdf, pp. 16, 15; p. 15 supplies definitions used on p. 16.}
\paragraph{Supplement: terminology.}
The notebook labels the likelihood derivative an efficient score:
\[
S(\theta)=\frac{d\log L(\theta)}{d\theta},
\]
where $L(\theta)$ is the likelihood. The empirical variance of $n$ sample scores is
\[
\frac1{n-1}\sum_i(S(\theta_i)-\overline{S(\theta)})^2.
\]
Efficient variance is the smallest attainable variance for a given problem, with a Cramer--Rao lower-bound interpretation.
\textit{Translator's note: the derivative shown is the ordinary likelihood score; calling it the efficient score requires the relevant model/orthogonalization context.}
\paragraph{Variance estimation.}
Define the oracle summand
\[
\psi_i^*=\mu_{(1)}(X_i)-\mu_{(0)}(X_i)
+\frac{W_i(Y_i-\mu_{(1)}(X_i))}{e(X_i)}
-\frac{(1-W_i)(Y_i-\mu_{(0)}(X_i))}{1-e(X_i)}.
\]
Its empirical variance converges to the efficient variance:
\[
\frac1{n-1}\sum_i(\psi_i^*-\widehat\tau^*)^2\longrightarrow_p V^*.
\]
The margin calls $\widehat\tau^*$ the non-cross-fitted oracle AIPW estimator. Replacing each nuisance function by the estimate trained outside $k(i)$ gives the cross-fitted summand $\widehat\psi_i$ and
\[
\widehat V_{\rm AIPW}=\frac1{n-1}\sum_i(\widehat\psi_i-\widehat\tau_{\rm AIPW})^2\longrightarrow_p V^*.
\]
The confidence interval drawn under a central normal probability region is
\[
\widehat\tau_{\rm AIPW}\ \pm\ \frac1{\sqrt n}\Phi^{-1}(1-\alpha/2)\sqrt{\widehat V_{\rm AIPW}}.
\]
\textit{Translator's note: the handwritten text says ``$\alpha$-level,'' while the formula and shaded sketch correspond to coverage $1-\alpha$.}
\section{Lecture 4: Estimating Treatment Heterogeneity}
\sourcepages{Lecture1-10.pdf, pp. 17--18.}
For observations $(X_i,Y_i,W_i)$ with $X_i\in\mathcal X$, we previously pursued the ATE $E[Y_i(1)-Y_i(0)]$. Heterogeneity analysis instead considers the conditional average treatment effect (CATE):
\[
\tau(x)=E[Y_i(1)-Y_i(0)\mid X_i=x].
\]
The blue note describes this as the average for a particular category of $X_i$.
\subsection{Regularization bias}
Under $\{Y_i(0),Y_i(1)\}\perp W_i\mid X_i$,
\[
\tau(x)=\mu_{(1)}(x)-\mu_{(0)}(x),\qquad
\mu_{(w)}(x)=E[Y_i\mid X_i=x,W_i=w].
\]
A direct CATE estimator is $\widehat\tau_T(x)=\widehat\mu_{(1)}(x)-\widehat\mu_{(0)}(x)$.
\paragraph{Personal concern.}
If there are too few observations for $\widehat\mu_{(1)}$, it may be heavily regularized (toward a straight line). If there are many observations for $\widehat\mu_{(0)}$, that fit may be accurate; then the difference $\widehat\tau_T(x)$ goes wrong because of the first fit.
\subsection{Semiparametric modeling}
\[
Y_i(w)=f(X_i)+w\tau(X_i)+\varepsilon_i(w),\qquad
P(W_i=1\mid X_i)=e(X_i),\qquad \tau(x)=\psi(x)\beta.
\]
$f$ is a function of $X_i$, and $\psi(x)$ is a nonparametric representation of $X$; $\beta$ is the treatment-effect coefficient. Under unconfoundedness, the Robinson transformation is
\[
Y_i-m(X_i)=(W_i-e(X_i))\psi(X_i)\beta+\varepsilon_i,
\quad m(x)=E[Y_i\mid X_i=x]=f(x)+e(x)\tau(x),
\]
with $\varepsilon_i=\varepsilon_i(W_i)$.
Design a residual-on-residual OLS regression using
\[
\widetilde Y_i^*=Y_i-m(X_i),\qquad
\widetilde Z_i^*=(W_i-e(X_i))\psi(X_i).
\]
Regress $\widetilde Y_i^*$ on $\widetilde Z_i^*$ to obtain $\widehat\beta_R^*$ (the notebook uses a coefficient symbol resembling $\zeta_R$). These transformations use only the observed data and the nuisance functions; the margin calls this data preprocessing. The notebook gives
\[
\sqrt n(\widehat\beta_R^*-\beta)\Rightarrow N(0,V_R),
\]
\[
V_R=\operatorname{Var}[\widetilde Z_i^*]^{-1}
\operatorname{Var}[\widetilde Z_i^*\widetilde Y_i^*]
\operatorname{Var}[\widetilde Z_i^*]^{-1}.
\]
Under homoskedasticity, $\operatorname{Var}(\varepsilon_i\mid X_i,W_i)=b^2$ is constant and $V_R$ is the semiparametrically efficient variance for $\beta$. Under heteroskedasticity, the notes say that $V_R$ still applies but is not the efficient variance.
\textit{Translator's note: the middle sandwich term above is reproduced literally from the scan; a regression sandwich calculation ordinarily uses the regression residual, rather than the unadjusted transformed outcome.}
\paragraph{Feasible cross-fitted version.}
The preceding regression is idealized because $m(X_i)$ and $e(X_i)$ are unknown. Fit nonparametric regressions $Y\sim X$ and $W\sim X$; omit the fold containing observation $i$ when constructing
\[
\widetilde Y_i=Y_i-\widehat m^{(-k(i))}(X_i),\qquad
\widetilde Z_i=(W_i-\widehat e^{(-k(i))}(X_i))\psi(X_i).
\]
Regress $\widetilde Y_i$ on $\widetilde Z_i$ to obtain $\widehat\beta_R$. If both nuisance RMSEs are $o_p(n^{-1/4})$, the notes conclude
\[
\sqrt n(\widehat\beta_R^*-\widehat\beta_R)\longrightarrow_p0.
\]
\subsection{A loss function for treatment heterogeneity}
\sourcepages{Lecture1-10.pdf, pp. 20, 19; p. 19 supplements the residual regression and the R-loss on p. 20.}
The previous estimator assumes $\tau(x)=\psi(x)\beta$. Often the relationship is not linear, so use machine learning. With
\[
\mu_{(w)}(x)=E[Y(w)\mid X=x],\quad
\varepsilon_i(w)=Y_i(w)-[\mu_{(0)}(X_i)+w\tau(X_i)],
\]
unconfoundedness yields $E[\varepsilon_i(w)\mid X_i,W_i]=0$. Here $\mu_{(0)}$ plays the role of the previous $f$. Again,
\[
Y_i-m(X_i)=(W_i-e(X_i))\tau(X_i)+\varepsilon_i,
\quad m(x)=\mu_{(0)}(x)+e(x)\tau(x),\quad \varepsilon_i=\varepsilon_i(W_i).
\]
Equivalently,
\[
\tau(\cdot)=\arg\min_{t}E[((Y_i-m(X_i))-(W_i-e(X_i))t(X_i))^2].
\]
If $m$ and $e$ are known, minimizing this gives $\tau(\cdot)$. The oracle penalized estimator is
\[
\widehat\tau_R^*(\cdot)=\arg\min_t\left\{
\frac1n\sum_i[(Y_i-m(X_i))-(W_i-e(X_i))t(X_i)]^2+\Lambda_n(t)\right\}.
\]
$\Lambda_n$ regularizes the complexity of the treatment-effect function, for example with lasso or ridge regression; the margin also mentions CNNs. Usually $m$ and $e$ are unknown (the notes parenthetically mention RCTs), so cross-fit them:
\begin{align*}
\widehat L_n(t)&=\frac1n\sum_i
[(Y_i-\widehat m^{(-k(i))}(X_i))-(W_i-\widehat e^{(-k(i))}(X_i))t(X_i)]^2,\\
\widehat\tau_R&=\arg\min_t\{\widehat L_n(t)+\Lambda_n(t)\}.
\end{align*}
For lasso, $\widehat\tau(x)=x\widehat\beta$, with
\[
\widehat\beta=\arg\min_\beta\left\{\frac1n\sum_i
[(Y_i-\widehat m^{(-k(i))}(X_i))-(W_i-\widehat e^{(-k(i))}(X_i))X_i\beta]^2
+\lambda\|\beta\|_1\right\}.
\]
\paragraph{Personal supplementary observation.}
With a constant treatment effect,
$Y_i(w)=f(X_i)+\tau w+\varepsilon_i$ and $E[\varepsilon_i\mid X_i,W_i]=0$, the residual estimator is
\[
\widehat\tau_R=\frac{\sum_i(Y_i-\widehat m^{(-k(i))}(X_i))(W_i-\widehat e^{(-k(i))}(X_i))}
{\sum_i(W_i-\widehat e^{(-k(i))}(X_i))^2}.
\tag{4.5 in the notes}
\]
The premise is that the nuisance components converge fast enough. Even if $\tau$ is not constant, this estimator converges to a nonnegatively weighted average of $\tau(x)$ (the annotation also writes a treated-minus-control mean comparison). When heterogeneity is small and overlap presents difficulties, use (4.5) as a replacement. The arrow connects this observation to the R-loss.

\section{Lecture 5: Regression Discontinuity Designs}
\sourcepages{Lecture1-10.pdf, pp. 22, 21; p. 21 explains why the earlier IPW approach does not apply.}
Introduce a cutoff using a running variable $Z_i\in\mathbb R$:
\[
W_i=\mathbf1_{\{Z_i\geq c\}},\qquad c\text{ is the threshold.}
\]
As $Z_i\to c$, observations near $c$ have similar characteristics and differ in whether they receive the intervention. Let $\mu_{(w)}(z)=E[Y_i(w)\mid Z_i=z]$. If the two conditional means are continuous, the CATE at $z=c$ is
\[
\tau_c=\mu_{(1)}(c)-\mu_{(0)}(c)
=\lim_{z\to c^+}E[Y_i\mid Z_i=z]-\lim_{z\to c^-}E[Y_i\mid Z_i=z].
\]
\paragraph{Personal supplement: why does ordinary IPW not apply?}
Earlier we assumed unconfoundedness $\{Y_i(0),Y_i(1)\}\perp W_i\mid Z_i$ and overlap
$\eta\leq P(W_i=1\mid Z_i)\leq1-\eta$. In RDD, unconfoundedness holds, but overlap fails because the treatment probability given $Z_i=z$ is either zero or one. The blue comment says this fails regardless of the nondegenerate choice of $\eta$. We must compare observations near $c$, using similarity rather than overlap.
Failure of overlap means $\tau_c$ cannot be estimated at the $\sqrt n$ rate. Its minimax error converges at a nonparametric rate depending on the smoothness of $\mu_{(w)}(z)$. For example, a uniformly bounded second derivative at $c$ allows an $n^{-2/5}$ estimation-error rate.
\subsection{Estimation via local linear regression}
Choose a small bandwidth $h_n\to0$ and a symmetric weighting function $K$. Fit the two conditional means locally. The coefficient on $W_i$ in
\[
\arg\min_{a,\tau,\beta_{(0)},\beta_{(1)}}\sum_i
K\left(\frac{|Z_i-c|}{h_n}\right)
[Y_i-a-\tau W_i-\beta_{(0)}(Z_i-c)_--\beta_{(1)}(Z_i-c)_+]^2
\]
is $\widehat\tau_c$. The intercept $a$ and slopes $\beta_{(w)}$ are nuisance parameters; the handwritten labels identify negative-side and positive-side parts. Common kernels are the window function $K(x)=\mathbf1_{\{|x|\leq1\}}$ and triangular kernel $K(x)=(1-|x|)_+$.
\subsection{Estimator properties and smoothness}
\sourcepages{Lecture1-10.pdf, pp. 23--24.}
Quantify continuity by assuming $\mu_{(w)}$ is twice differentiable with a uniformly bounded second derivative:
\[
\left|\frac{d^2}{dz^2}\mu_{(w)}(z)\right|\leq B,
\qquad z\in\mathbb R,\quad w\in\{0,1\}.
\]
If smoothness is weaker, with only Lipschitz continuity (the note says bounded but not differentiable), local averaging may work better because there is no clear linear trend locally. If smoothness is stronger, with bounded derivatives of order $k\geq3$, higher-order local regression may improve the convergence rate.
Taylor-expand at $c$:
\[
\mu_{(w)}(z)=a_{(w)}+\beta_{(w)}(z-c)+\tfrac12\rho_{(w)}(z-c),
\qquad |\rho_{(w)}(u)|\leq Bu^2.
\]
The last term is the remainder, and $\tau_c=a_{(1)}-a_{(0)}$. Fit each side separately:
\[
(\widehat a_{(1)},\widehat\beta_{(1)})=arg\min_{a,\beta}
\sum_{Z_i\geq c}K\left(\frac{|Z_i-c|}{h_n}\right)[Y_i-a-\beta(Z_i-c)]^2,
\]
and analogously on the control side; then $\widehat\tau_c=\widehat a_{(1)}-\widehat a_{(0)}$.
\paragraph{Closed-form weights for the window kernel.}
On $c\leq Z_i\leq c+h_n$, the scan writes
\[
\widehat a_{(1)}=\sum_i\gamma_iY_i,\qquad
\gamma_i=\frac{\widehat E_{(1)}[(Z_i-c)^2]-\widehat E_{(1)}[Z_i-c](Z_i-c)}
{\widehat E_{(1)}[(Z_i-c)^2]-\widehat E_{(1)}[Z_i-c]^2},
\]
where $\widehat E_{(1)}[Z_i-c]$ is the average of $Z_i-c$ within the bandwidth. It also records
\[
\sum_i\gamma_i=1,\qquad \sum_i\gamma_i(Z_i-c)=0.
\]
\textit{Translator's note: the displayed weight in the scan is missing division by the number of observations in the bandwidth. That normalization is required for the subsequent sum-to-one identity. All sums in this paragraph run over the bandwidth.}
The decomposition on the page is
\[
\widehat a_{(1)}=a_{(1)}+\underbrace{\sum_i\gamma_i\rho_{(1)}(Z_i-c)}_{\text{curvature bias}}
+\underbrace{\sum_i\gamma_i(Y_i-\mu_{(1)}(Z_i))}_{\text{sampling noise}}.
\]
The control intercept has an analogous decomposition. The red arrows say that both terms need bounds. Curvature bias is bounded at order $Bh_n^2$, where $B$ bounds the second derivative and $h_n$ is the bandwidth. Sampling noise has mean zero; if $\operatorname{Var}[Y_i\mid Z_i]\leq b^2$, its variance is bounded by $b^2\sum_i\gamma_i^2$.
If $Z_i$ has a continuous nonzero density $f$ near $c$,
\[
b^2\sum_i\gamma_i^2\approx\frac{4b^2}{|\{i:c\leq Z_i\leq c+h_n\}|}
\approx\frac{4b^2}{f(c)nh_n}.
\]
Thus squared bias is of order $h_n^4$ and variance is of order $(nh_n)^{-1}$, giving
\[
\widehat\tau_c=\tau_c+O_p(n^{-2/5}),\qquad h_n\asymp n^{-1/5}.
\]
The red note identifies $n^{-2/5}$ as the estimation error and adds that with $k$-th order smoothness, $h_n\asymp n^{-1/(2k+1)}$.
\textit{Translator's note: the remainder decomposition drops the factor $1/2$ used in the preceding Taylor line; this does not change the rate argument.}
\subsection{Noise in the running variable}
\sourcepages{Lecture1-10.pdf, p. 25.}
The blue margin notes that the running variable may contain noise. Introduce a latent $U_i\sim G$ and two conditions:
\begin{enumerate}
\item A noisy running variable: the scan writes $Z_i\mid U_i\sim N(0,\nu^2)$, with $\nu>0$.
\item Unconfoundedness of the running variable: $\{Y_i(0),Y_i(1)\}\perp Z_i\mid U_i$.
\end{enumerate}
With $Y_i=Y_i(W_i)$ and $\{Y_i(0),Y_i(1)\}\perp W_i\mid U_i$,
\[
e(u)=P(W_i=1\mid U_i=u)=P(Z_i\geq c\mid U_i=u)
=1-\Phi\left(\frac{c-u}{\nu}\right),
\]
\[
\alpha_{(w)}(u)=E[Y_i(w)\mid U_i=u],\qquad
\tau(u)=E[Y_i(1)-Y_i(0)\mid U_i=u].
\]
The annotation defines the standard normal CDF,
$\Phi(x)=(2\pi)^{-1/2}\int_{-\infty}^x e^{-t^2/2}\,dt$.
\textit{Translator's note: the ensuing propensity and density formulas require $Z_i\mid U_i=u\sim N(u,\nu^2)$, or equivalently $Z_i-U_i\mid U_i\sim N(0,\nu^2)$. The conditional normal line above is preserved as written.}
$U_i$ is unobserved, so we cannot directly use its propensity score. Consider a deconvolution-type estimator:
\[
\widehat\tau_\gamma=\frac1n\sum_{Z_i\geq c}\gamma_+(Z_i)Y_i
-\frac1n\sum_{Z_i<c}\gamma_-(Z_i)Y_i.
\]
Choose the weights to satisfy
\[
E[\mathbf1_{\{Z_i\geq c\}}\gamma_+(Z_i)]
=E[\mathbf1_{\{Z_i<c\}}\gamma_-(Z_i)]=1.
\]
Then
\[
E[\widehat\tau_\gamma]
=\underbrace{\int h_+(u)\tau(u)\,dG(u)}_{\text{weighted treatment effect}}
+\underbrace{\int(h_+(u)-h_-(u))\alpha_{(0)}(u)\,dG(u)}_{\text{confounding bias}},
\]
where the weighting functions are
\[
h_+(u)=\int_c^\infty\gamma_+(z)\varphi_\nu(z-u)\,dz,\qquad
h_-(u)=\int_{-\infty}^c\gamma_-(z)\varphi_\nu(z-u)\,dz,
\]
\[
\varphi_\nu(z-u)=\frac1{\sqrt{2\pi\nu^2}}\exp\left[-\frac12\left(\frac{z-u}{\nu}\right)^2\right].
\]

\section{Lecture 6: Finite-Sample Inference in RDDs}
\sourcepages{Lecture1-10.pdf, pp. 27, 26, 29, 28; p. 26 continues p. 27, and p. 28 explains the weight constraints on p. 29.}
Finite-sample RDDs allow a discontinuous/discrete running variable or a more complex cutoff rule. The local-linear regression solution (6.2) can be written as a linear estimator,
\[
\widehat\tau_c(\gamma)=\sum_{i=1}^n\gamma_iY_i.\tag{6.3 in the notes}
\]
The weights depend only on $Z_i$. If
\[
Y_i(w)=\mu_{(w)}(Z_i)+\varepsilon_i(w),\qquad
\varepsilon_i(w)\mid Z_i\sim N(0,b^2),
\]
then, conditional on $Z_1,\ldots,Z_n$,
\[
\widehat\tau_c(\gamma)\sim N(\widehat\tau_c^*(\gamma),b^2\|\gamma\|_2^2),
\qquad \widehat\tau_c^*(\gamma)=\sum_i\gamma_i\mu_{(W_i)}(Z_i).
\]
Here $W_i=\mathbf1_{\{Z_i\geq c\}}$ and $\|\gamma\|_2=(\gamma_1^2+\cdots+\gamma_n^2)^{1/2}$ is the Euclidean norm.
\subsection{Minimax linear estimation}
Minimax means minimizing the worst-case MSE. For a linear estimator, the conditional variance is directly accessible as $b^2\|\gamma\|_2^2$ under homoskedasticity. Its bias depends on unknown conditional means:
\[
\operatorname{Bias}(\widehat\tau_c(\gamma)\mid Z_1,\ldots,Z_n)
=\sum_i\gamma_i\mu_{(W_i)}(Z_i)-[\mu_{(1)}(c)-\mu_{(0)}(c)].
\]
A smoothness restriction $|\mu_{(w)}''(z)|\leq B$ bounds the bias by
\[
I_B(\gamma)=\sup_{|\mu_{(w)}''|\leq B}
\left\{\sum_i\gamma_i\mu_{(W_i)}(Z_i)-[\mu_{(1)}(c)-\mu_{(0)}(c)]\right\}.
\]
The blue note says the supremum does not need an absolute value because the second-derivative class is symmetric. Hence
\[
|\operatorname{Bias}(\widehat\tau_c(\gamma)\mid Z_1,\ldots,Z_n)|\leq I_B(\gamma),
\]
\[
\operatorname{MSE}(\widehat\tau_c(\gamma)\mid Z_1,\ldots,Z_n)
\leq b^2\|\gamma\|_2^2+I_B^2(\gamma).
\]
The red annotation emphasizes ``MSE = variance + squared bias'' and that this bound does not depend on the unknown $\mu$. The left-page continuation defines
\[
\gamma^B=\arg\min_\gamma\{b^2\|\gamma\|_2^2+I_B^2(\gamma)\},\qquad
\widehat\tau_c(\gamma^B)=\sum_i\gamma_i^B Y_i.
\tag{6.8 in the notes}
\]
The blue comparison says that, at a fixed bias bound $I_B(\gamma)$, this yields smaller MSE than local-linear estimation.
\subsection{Deriving the minimax linear weights}
Decompose
\[
\mu_{(w)}(z)=a_{(w)}+\beta_{(w)}(z-c)+\rho_{(w)}(z),
\]
where $\rho_{(w)}(c)=\rho'_{(w)}(c)=0$ and the second derivative is bounded by $B$. The intercepts and slopes are unrestricted, so $I_B(\gamma)$ can be infinite unless
\[
\sum_i\gamma_iW_i=1,\qquad \sum_i\gamma_i=0,\qquad
\sum_i\gamma_i(Z_i-c)_+=0,\qquad \sum_i\gamma_i(Z_i-c)_-=0.
\]
After these constraints, only the remainders need control. Rewrite (6.8) as
\[
(\gamma^B,t)=\arg\min_{\gamma,t}\{b^2\|\gamma\|_2^2+B^2t^2\},
\]
subject to
\[
\sum_i\gamma_iW_i\rho_{(1)}(Z_i)+\sum_i\gamma_i(1-W_i)\rho_{(0)}(Z_i)\leq t
\]
for every $\rho_{(w)}$ with $\rho_{(w)}(c)=\rho'_{(w)}(c)=0$ and $|\rho''_{(w)}(z)|\leq1$, together with
\[
\sum_i\gamma_iW_i=1,\quad\sum_i\gamma_i=0,\quad
\sum_i\gamma_iW_i(Z_i-c)=0,\quad
\sum_i\gamma_i(1-W_i)(Z_i-c)=0.
\tag{6.10 in the notes}
\]
The blue instruction is to solve the corresponding dual problem.
\paragraph{Personal supplement: why do the weights sum to zero?}
The final estimate is a difference. Split its sum into the treated and control parts:
\[
\widehat\tau_c(\gamma)=\sum_i\gamma_i^{(1)}Y_i+\sum_i\gamma_i^{(2)}Y_i.
\]
The treated weights sum to $1$ and the control weights sum to $-1$, so
\[
\sum_i\gamma_i=1+(-1)=0,\quad
\sum_i\gamma_i^{(1)}=\sum_i\gamma_iW_i=1,\quad
\sum_i\gamma_i^{(2)}=\sum_i\gamma_i(1-W_i)=-1.
\]
The sketch shows a weight profile with a negative spike just to the left of the threshold and a positive spike just to the right; farther away, the weights change sign and then approach zero. The red labels on the original say $Y=1$ and $Y=0$, while the equations below them use treatment indicators $W$.
\subsection{Inference with linear estimators}
\sourcepages{Lecture1-10.pdf, pp. 30--31.}
Construct a confidence interval for $\widehat\tau_c(\gamma^B)$. Let
\[
\mathrm{err}=\widehat\tau_c-\tau_c,\qquad
\mathrm{err}\mid Z_1,\ldots,Z_n\sim N(\mathrm{bias},b^2\|\gamma^B\|_2^2),
\]
where $|\mathrm{bias}|\leq Bt$ from (6.10). By unimodality of the Gaussian distribution, for $S\sim N(0,1)$,
\[
P(|\mathrm{err}|\geq\zeta)\leq P(|Bt+b\|\gamma^B\|_2S|\geq\zeta).
\]
A confidence interval with conditional coverage at least $1-\alpha$ is
\[
I_\alpha=(\widehat\tau_c(\gamma^B)-\zeta_\alpha^B,
\widehat\tau_c(\gamma^B)+\zeta_\alpha^B),
\]
\[
\zeta_\alpha^B=\inf\{\zeta:P(|Bt+b\|\gamma^B\|_2S|>\zeta)\leq\alpha\}.
\tag{6.13 in the notes}
\]
The blue comment stresses that conditioning on $Z_i$ makes this work for running variables with any distribution, including nonstandard RDD settings.
\paragraph{Discrete running variables.}
The restriction $|\mu_{(w)}''(z)|\leq B$ alone need not identify $\tau_c$, but (6.13) remains valid. For continuous $Z_i$, the interval shrinks at rate $n^{-2/5}$; for discrete $Z_i$, its length may not tend to zero. When $c$ is outside the support of the running variable, one needs an interpolating function representing $\mu_{(w)}(z)$ and satisfying the second-derivative bound.
\paragraph{Multivariate running variables.}
Examples include two- or three-dimensional regions and geographic boundaries. A scalar-cutoff RDD cannot be used directly; some reduce the problem to one-dimensional distance, which the note calls less stable. The minimax method can instead be extended directly to multiple dimensions.
\subsection{Beyond homoskedasticity}
When the errors $\varepsilon_i=Y_i-\mu_{(W_i)}(Z_i)$ are no longer homoskedastic Gaussian with variance $b^2$, the exact normal statement (6.11) fails. Use a central limit theorem to justify
\[
\widehat\tau_c(\gamma)\mid Z_1,\ldots,Z_n
\ \approx\ N\left(\widehat\tau_c^*(\gamma),
\sum_i\gamma_i^2\operatorname{Var}[Y_i\mid Z_i,W_i]\right).
\tag{6.14 in the notes}
\]
With a valid Gaussian approximation, confidence intervals can still be obtained. Estimate the conditional variance by
\[
\widehat V_n=\sum_i\gamma_i^2(Y_i-\widehat\mu_{(W_i)}(Z_i))^2,
\]
where the fitted mean comes from local-linear regression. The blue note says that if this mean model is wrong, residuals are inflated, making the result too conservative. The minimax criterion (6.8) was for homoskedasticity; under heteroskedasticity, introduce individual variances $b_i^2$ into (6.10), the problem used to construct the dual.
\section{Lecture 7: Balancing Estimators}
\sourcepages{Lecture1-10.pdf, pp. 32--33.}
The heading links balance to the principle behind estimating $\widehat e(x)$. Two aims are noted: improving the propensity and regression functions, and estimating ATEs with high-dimensional covariates. IPW and AIPW use propensity scores; this lecture uses convex optimization to derive how to estimate $e(x)$.
\subsection{Review of IPW}
For IID observations $(X_i,Y_i,W_i)\in\mathcal X\times\mathbb R\times\{0,1\}$, the goal is $\tau=\mu_{(1)}-\mu_{(0)}$, where $\mu_{(w)}=E[Y_i(w)]$. Assume
\[
\{Y_i(0),Y_i(1)\}\perp W_i\mid X_i,\qquad
0<\eta\leq e(X_i)\leq1-\eta<1,
\]
with $e(x)=P(W_i=1\mid X_i=x)$ and positive constant $\eta$. Estimate the treated mean using
\[
\widehat\mu_{\rm IPW}(1)=\frac1n\sum_i\frac{W_iY_i}{\widehat e(X_i)}.
\tag{7.3 in the notes}
\]
The notebook labels the following consistency argument ``brief, not important'' but records three steps:
\begin{enumerate}
\item Population balance: under unconfoundedness, $\mu_{(1)}=E[W_iY_i/e(X_i)]$.
\item Oracle estimator: with the true $e(x)$ it is unbiased; overlap gives finite variance.
\item Feasible approximation: if the estimated propensities also satisfy overlap, the Cauchy--Schwarz inequality gives
\begin{align*}
|\widehat\mu_{\rm IPW}(1)-\widehat\mu^*_{\rm IPW}(1)|
&\leq\left[\frac1n\sum_i\left(\frac1{\widehat e(X_i)}-\frac1{e(X_i)}\right)^2\right]^{1/2}
\left[\frac1n\sum_i(W_iY_i)^2\right]^{1/2}\\
&\leq\frac1{\eta^2}\left[\frac1n\sum_i(\widehat e(X_i)-e(X_i))^2\right]^{1/2}
\left[\frac1n\sum_i(W_iY_i)^2\right]^{1/2}.
\end{align*}
The margin recalls $|a\cdot b|^2\leq\|a\|^2\|b\|^2$. It adds that the analogous AIPW error is negligible relative to $1/\sqrt n$ as $n\to\infty$.
\end{enumerate}
\subsection{Population versus sample balance}
Expand the conditional mean in a predefined basis:
\[
\mu_{(w)}(x)=\sum_{j=1}^\infty\beta_{j(w)}\psi_j(x).
\]
The red note gives polynomial and Fourier bases as examples. Under reasonable regularity conditions,
\[
\mu_{(w)}=\sum_{j=1}^\infty\beta_{j(w)}E[\psi_j(X_i)].
\]
The arrow says to take expectations of both sides; a red question asks whether these assumptions may have a problem.
With unconfoundedness, $Y_i=\mu_{(W_i)}(X_i)+\varepsilon_i$ and $E[\varepsilon_i\mid X_i,W_i]=0$. The population IPW argument is
\begin{align*}
E\left[\frac{W_iY_i}{e(X_i)}\right]
&=E\left[\frac{W_i}{e(X_i)}\sum_j\beta_{j(1)}\psi_j(X_i)\right]\\
&=\sum_j\beta_{j(1)}E\left[\frac{W_i\psi_j(X_i)}{e(X_i)}\right]
=\sum_j\beta_{j(1)}E[\psi_j(X_i)]=\mu_{(1)}.
\end{align*}
The red explanation identifies $E[W_i\mid X_i]=e(X_i)$ as the core of the argument. Thus choose weights so that, in the sample,
\[
\frac1n\sum_i\frac{W_i\psi_j(X_i)}{\widehat e(X_i)}
\approx\frac1n\sum_i\psi_j(X_i),\qquad j=1,2,\ldots.
\tag{7.8 in the notes}
\]
\subsection{A balancing loss for propensity estimation}
\sourcepages{Lecture1-10.pdf, pp. 35, 34; p. 34 supplies the norm reminders used in the high-dimensional bounds.}
Suppose the outcome model is linear, $\mu_{(w)}(x)=x\beta_{(w)}$, and the propensity model is logistic,
\[
e(x)=\frac1{1+\exp(-x\theta)}.
\]
The blue annotation says that sample balance then involves only the covariates $X_i$, and the targeted balance concerns means. Equation (7.8) becomes
\[
\frac1n\sum_i(1+e^{-X_i\widehat\theta})W_iX_i=\frac1n\sum_iX_i.
\tag{7.9 in the notes}
\]
This is the KKT condition of the convex problem
\[
\widehat\theta=\arg\min_\theta\frac1n\sum_i\ell_\theta(X_i,Y_i,W_i),\qquad
\ell_\theta(X_i,Y_i,W_i)=W_ie^{-X_i\theta}+(1-W_i)X_i\theta.
\tag{7.10 in the notes}
\]
It can therefore be solved by Newton descent. The handwritten text describes a unique solution and says this serves as the loss function. The margin notes that estimating both $\mu_{(0)}$ and $\mu_{(1)}$ requires two propensity models.
\subsection{ATE estimation with high-dimensional confounders}
Suppose $X_i\in\mathbb R^p$, possibly $p\gg n$, and $\mu_{(w)}(x)=x\beta_{(w)}$. In high dimensions, it is difficult to estimate $e(x)$ adequately; exact weights are not attainable, so seek approximate balance. By Holder's inequality,
\begin{align*}
&\frac1n\sum_i\frac{W_i\mu_{(1)}(X_i)}{\widehat e(X_i)}-\frac1n\sum_i\mu_{(1)}(X_i)\\
&\quad=\left(\frac1n\sum_i\frac{W_iX_i}{\widehat e(X_i)}-\frac1n\sum_iX_i\right)\beta_{(1)}\\
&\quad\leq\left\|\frac1n\sum_i\frac{W_iX_i}{\widehat e(X_i)}-\frac1n\sum_iX_i\right\|_\infty\|\beta_{(1)}\|_1.
\end{align*}
The first norm is labeled imbalance. The handwritten reminder states Holder's bound
$\sum_i|a_ib_i|\leq(\sum_i|a_i|^p)^{1/p}(\sum_i|b_i|^q)^{1/q}$.
\paragraph{Personal interpretation and questions.}
Good inverse-probability weights make the worst case of imbalance as good as possible. Two problems remain: (1) how to parameterize $\widehat e(X_i)$ to obtain sup-norm approximate balance; (2) the bound is meaningful only if $\|\beta_{(1)}\|_1$ is bounded. High-dimensional statistics usually assumes a sparse coefficient vector rather than constraining its $1$-norm directly.
\paragraph{Supplement: norm reminders.}
The left page lists the ``0-norm'' as $\sum_iX_i^0$, the $1$-norm as $\sum_i|X_i|$, the $2$-norm as $(\sum_iX_i^2)^{1/2}$, and the infinity norm as $\max_i|X_i|$. It gives the general expression $\|x\|_p=(\sum_i x_i^p)^{1/p}$.
\textit{Translator's note: these last formulas retain the handwritten shorthand. The standard $0$-norm counts nonzero entries, and the standard general $p$-norm uses $|x_i|^p$.}
\subsection{Approximate balancing weights and regression augmentation}
\sourcepages{Lecture1-10.pdf, pp. 36--37.}
For a well-chosen tuning parameter $\zeta$, solve
\[
\widehat\gamma=\arg\min_{\gamma_i\geq1}
\left\{\frac1{n^2}\|\gamma\|_2^2+
\zeta\left\|\frac1n\sum_i(\gamma_iW_i-1)X_i\right\|_\infty^2\right\}.
\tag{7.12 in the notes}
\]
The red labels identify the $L_2$ norm, imbalance, and tuning weight. Then use $1/\widehat e(X_i)=\widehat\gamma_i$ as inverse-propensity weights. Sub-Gaussian concentration inequalities verify, for oracle weights, bounds of the form
\[
\frac1{n^2}\|1/e(X_i)\|_2^2\leq\frac{\eta^{-2}}n,
\qquad
\left\|\frac1n\sum_i(\gamma_iW_i-1)X_i\right\|_\infty^2
=O_p\left(\frac{\eta^{-2}\log p}{n}\right).
\]
The notes say $\widehat e(X_i)$ need not estimate $e(X_i)$ well; with a suitable $\zeta$, these inverse weights nonetheless satisfy a bias bound.
Next augment the weighted estimator with high-dimensional regression to remove dependence on $\|\beta_{(1)}\|_1$. The augmented balancing estimator is
\[
\widehat\mu_{\rm AB}(1)=\frac1n\sum_i
\left[\widehat\mu_{(1)}(X_i)+\widehat\gamma_iW_i(Y_i-\widehat\mu_{(1)}(X_i))\right],
\quad\widehat\mu_{(1)}(X_i)=X_i\widehat\beta_{(1)}.
\tag{7.14 in the notes}
\]
Its decomposition is
\begin{align*}
\widehat\mu_{\rm AB}(1)={}&\underbrace{\frac1n\sum_iX_i\beta_{(1)}}_{\text{mean of }\mu_{(1)}(X_i)}
+\underbrace{\frac1n\sum_i\widehat\gamma_iW_i(Y_i-X_i\beta_{(1)})}_{\text{mean-zero noise}}\\
&+\left[\frac1n\sum_i(1-\widehat\gamma_iW_i)X_i\right](\widehat\beta_{(1)}-\beta_{(1)}).
\end{align*}
The last term is bounded in magnitude by
\[
\left\|\frac1n\sum_i(1-\widehat\gamma_iW_i)X_i\right\|_\infty
\|\widehat\beta_{(1)}-\beta_{(1)}\|_1.
\]
The red intuition is that the weights only need to correct the estimation error $\beta_{(1)}-\widehat\beta_{(1)}$, rather than the whole fitted coefficient vector. High-dimensional regression supplies a bound on this error: under sparsity $\|\beta_{(1)}\|_0\leq k$ and a restricted-eigenvalue condition, lasso achieves
\[
\|\widehat\beta_{(1)}-\beta_{(1)}\|_1
=O_p\left(k\sqrt{\frac{\log p}{n}}\right).
\]
\paragraph{Putting the steps together.}
Suppose the design satisfies the conditions for this bound, $\beta_{(1)}$ is $k$-sparse, and $k\ll\sqrt n/\log p$.
\begin{enumerate}
\item Fit lasso on the treated units; the coefficient $1$-norm error is $o_p(1/\sqrt{\log p})$.
\item Fit $\widehat\gamma$ using (7.12), giving infinity-norm imbalance $O_p(\sqrt{\log(p)/n})$.
\item Use (7.14), and insert these bounds into the decomposition:
\[
\widehat\mu_{\rm AB}(1)-\frac1n\sum_iX_i\beta_{(1)}
=\frac1n\sum_i\widehat\gamma_iW_i(Y_i-X_i\beta_{(1)})+o_p(n^{-1/2}).
\tag{7.17 in the notes}
\]
The noise term is $O_p(n^{-1/2})$.
\end{enumerate}
A central limit theorem also holds:
\[
\frac{\sum_i\widehat\gamma_iW_i(Y_i-X_i\beta_{(1)})}
{\sqrt{\sum_i\widehat\gamma_i^2W_i^2(Y_i-X_i\beta_{(1)})^2}}
\Rightarrow N(0,1).
\]
\paragraph{Personal caveat.}
This is weaker than expected: it only gives a confidence interval for the sample average
$\overline\mu_{(1)}=n^{-1}\sum_iX_i\beta_{(1)}$, not for $\mu_{(1)}$ itself. To obtain a result for the population mean, one needs to establish convergence of the estimated weights $\widehat\gamma_i$.

\section{Lecture 8: Methods for Panel Data}
\sourcepages{Lecture1-10.pdf, pp. 39, 38; p. 38 supplies the panel layout and within-unit variance supplement.}
Panel (longitudinal) data observe units $i=1,\ldots,n$ at times $t=1,\ldots,T$. The left-page sketch is a matrix with units down the rows and time points $t_1,t_2,\ldots,T$ across the columns.
\subsection{A constant treatment-effect model}
Outcomes are $Y_{it}\in\mathbb R$ and treatment assignments are $W_{it}\in\{0,1\}$. The basic model is
\[
Y_{it}=Y_{it}(0)+W_{it}\tau,\qquad i=1,\ldots,n,\quad t=1,\ldots,T.
\tag{8.1 in the notes}
\]
This very basic model leaves out treatment heterogeneity and treatment dynamics. The notes explain that the treatment effect is the same for all units and all times; a unit's outcome at $t$ depends only on its treatment at $t$, not on previous treatment. The red annotations emphasize that $\tau$ is unrelated to $i,t$ and that past times do not enter.
\subsection{The two-way model}
\[
Y_{it}=\alpha_i+\beta_t+W_{it}\tau+\varepsilon_{it},
\qquad E[\varepsilon\mid\alpha,\beta,W]=0.
\tag{8.2 in the notes}
\]
$\alpha_i$ and $\beta_t$ are unit and time shifts. The red comment calls the condition strict exogeneity of treatment, stronger than unconfoundedness. For example, with two periods, let one group have treatment path $(0,0)$ and another $(0,1)$. The closed-form difference-in-differences estimate is
\[
\widehat\tau=\frac1{|\{i:W_{i2}=1\}|}\sum_{W_{i2}=1}(Y_{i2}-Y_{i1})
-\frac1{|\{i:W_{i2}=0\}|}\sum_{W_{i2}=0}(Y_{i2}-Y_{i1}).
\]
\paragraph{Comments in the notes.}
A limitation is that with large $n,T$ and general treatment-assignment distributions, model (8.2) may fail. Analogously, OLS with more than two periods is sometimes also called DID, although the note says there is no closed-form OLS solution. The listed advantage is the intuitive parallel-trends requirement (units differ only in $\alpha_i$); without parallel trends, it cannot be used.
\textit{Translator's note: the phrase about there being no closed-form OLS solution is retained as written; general OLS does have a matrix expression. The handwritten point appears to concern the absence of the simple two-period DID expression.}
\paragraph{Supplement: modeling the noise and estimating variance.}
Treat the row vectors $\varepsilon_i=(\varepsilon_{i1},\ldots,\varepsilon_{iT})$ as IID across units, with $\operatorname{Var}[\varepsilon_i]=\Sigma$. Treating every $\varepsilon_{it}$ as independent is hard to justify and unreasonable. For $T=2$, estimate the variance of the differences directly:
\[
\widehat{\operatorname{Var}}[\widehat\tau]
=\frac{\widehat{\operatorname{Var}}[Y_{i2}-Y_{i1}\mid W_{i2}=1]}{|\{i:W_{i2}=1\}|}
+\frac{\widehat{\operatorname{Var}}[Y_{i2}-Y_{i1}\mid W_{i2}=0]}{|\{i:W_{i2}=0\}|}.
\tag{8.4 in the notes}
\]
This remains robust when errors within a row/unit are correlated.
% Continuation: stats1-tail.tex translates original pp.40--58, beginning with interactive panel models.
